\section{Transcripciones de entrevistas}
\label{anexo:transcripciones}

Se presentan las cinco transcripciones recuperadas del levantamiento del SGRA~\cite{SGRA-documentacion-original}. Se conservan las intervenciones y las marcas de tiempo disponibles. Los identificadores coinciden con el inventario del Anexo~\ref{anexo:evidencias}.

\begingroup
\small
\setstretch{1.1}

\subsection{ENT-01 — Jessica Ponce, coordinación}
\label{transcripcion:ent01}

Participantes: Bryan Lombeida, Mario Zambrano, Belinda Toaquiza, Melanie Muñoz (equipo de desarrollo). Modalidad: presencial. Duración aproximada: 15–20 minutos.\par\smallskip

\noindent\textbf{Bryan (0:} 06): Muy buenas tardes. Estamos aquí con mi persona y mis compañeros Mario, Belinda y Melanie. En esta ocasión queremos realizar una entrevista para conocer funcionalidades, requisitos o necesidades para el sistema que estamos desarrollando.\par\smallskip

\noindent\textbf{Mario (0:} 19): El sistema está enfocado en la gestión docente cuando un estudiante requiere tutoría.\par\smallskip

\noindent\textbf{Belinda (0:} 23): El propósito de esta entrevista es conocer su perspectiva y experiencia como coordinadora de la carrera de Software.\par\smallskip

\noindent\textbf{Melanie (0:} 29): La entrevista tomará aproximadamente 15 a 20 minutos. Ya con su consentimiento para grabar, comenzamos.\par\smallskip

Preguntas y respuestas:\par\smallskip

\noindent\textbf{Bryan (0:} 44): ¿Podría compartirnos un poco de información acerca de su rol como coordinadora de la carrera de Ingeniería en Software en la UTEQ?\par\smallskip

\noindent\textbf{Coordinadora (0:} 51): El rol de una coordinadora es coordinar los procesos académicos y administrativos de la carrera. Me encargo de la planificación académica, el seguimiento del rendimiento estudiantil y la supervisión general de las actividades académicas.\par\smallskip

\noindent\textbf{Mario (1:} 07): ¿En qué medida ha participado o supervisado tutorías académicas, ya sea en formato presencial o en línea?\par\smallskip

\noindent\textbf{Coordinadora (1:} 16): He llevado el proceso de tutorías a nivel general, pidiendo planificaciones a los tutores, revisando su cumplimiento, evaluando el rendimiento académico y diseñando estrategias de mejora. Sin embargo, el seguimiento directo lo realizan los coordinadores de área.\par\smallskip

\noindent\textbf{Belinda (1:} 59): ¿Cómo describiría su experiencia en el uso del SGA para fines académicos o administrativos?\par\smallskip

\noindent\textbf{Coordinadora (2:} 06): El SGA es muy útil. Allí tengo acceso a calificaciones, asistencia y entrega de tareas de los estudiantes. Esto permite identificar dificultades y coordinar estrategias de mejora con los docentes. También gestionamos ayudantías de cátedra desde allí, aunque no siempre tienen continuidad, ya que depende tanto de docentes como de estudiantes.\par\smallskip

\noindent\textbf{Melanie (3:} 18): ¿Qué métodos se emplean actualmente en la carrera para gestionar tutorías entre estudiantes y docentes?\par\smallskip

\noindent\textbf{Coordinadora (3:} 24): Generalmente los docentes utilizan sus horas de docencia para atender a estudiantes, aunque esas horas suelen ser limitadas porque manejan varias asignaturas. Además, se usan ayudantías de cátedra. Algunos docentes recurren a WhatsApp, envío de videos o clases online fuera de horario, como alternativas para apoyar a los estudiantes.\par\smallskip

\noindent\textbf{Bryan (5:} 05): ¿Existen políticas institucionales sobre la frecuencia de las tutorías?\par\smallskip

\noindent\textbf{Coordinadora (5:} 12): No. A los docentes se les asignan horas de preparación de clases que pueden usar también para atender estudiantes, pero no hay una regulación específica sobre tutorías individuales o grupales.\par\smallskip

\noindent\textbf{Mario (6:} 00): ¿Cómo debería un sistema permitir a los docentes registrar y actualizar su disponibilidad para tutorías?\par\smallskip

\noindent\textbf{Coordinadora (6:} 08): Podría hacerse configurando horarios, partiendo del distributivo docente. Que cada docente indique qué horas puede destinar semanalmente a tutorías.\par\smallskip

\noindent\textbf{Belinda (6:} 36): ¿Y cómo garantizar que los docentes respondan a las solicitudes en un plazo razonable?\par\smallskip

\noindent\textbf{Coordinadora (6:} 44): Lo más efectivo sería notificaciones instantáneas. Un docente no siempre revisará un sistema web, pero sí un mensaje de WhatsApp o una notificación en una app. Eso permitiría una respuesta más rápida.\par\smallskip

\noindent\textbf{Melanie (7:} 42): ¿Cómo facilitar la organización de tutorías grupales cuando varios estudiantes comparten dudas similares?\par\smallskip

\noindent\textbf{Coordinadora (7:} 54): El docente debería tener la opción de aceptar solicitudes grupales. El estudiante podría decidir si quiere una tutoría grupal o individual, según el tema y su comodidad. El sistema debería permitir configurar estas modalidades.\par\smallskip

\noindent\textbf{Bryan (9:} 39): ¿Podría el sistema limitar el número de tutorías por estudiante?\par\smallskip

\noindent\textbf{Coordinadora (9:} 43): Eso dependería de la disponibilidad del docente. En lo presencial no hay problema porque el aula permite varios estudiantes. En línea también hay más flexibilidad. Sin embargo, hay temas que requieren atención personalizada, como cálculo diferencial. En esos casos el docente puede dividir en grupos más pequeños según su criterio y disponibilidad.\par\smallskip

\noindent\textbf{Mario (11:} 23): ¿Qué reglas o controles considera necesarios en un sistema de tutorías?\par\smallskip

\noindent\textbf{Coordinadora (11:} 28): Debe estar vinculado a las asignaturas activas del semestre. Por ejemplo, si un docente no está dictando una materia ese semestre, no tendría sentido que ofrezca tutorías en ella. Esto evita confusiones y asegura que los estudiantes reciban apoyo del docente que está trabajando directamente con el curso.\par\smallskip

\noindent\textbf{Belinda (13:} 18): ¿Algún aspecto adicional que desee compartir acerca de su experiencia o expectativas para el sistema?\par\smallskip

\noindent\textbf{Coordinadora (13:} 24): Sería muy importante que los estudiantes también reciban notificaciones de sus tutorías, recordatorios del lugar, hora y modalidad (presencial o en línea). Además, el sistema debería ser flexible para permitir distintas plataformas de conexión como Zoom, Meet o Teams, según la preferencia de cada docente.\par\smallskip

\noindent\textbf{Melanie (14:} 37): Muchas gracias por su tiempo y sus aportes.\par\smallskip

\noindent\textbf{Coordinadora (14:} 41): Con gusto. Gracias a ustedes.\par\smallskip

\subsection{ENT-02 — Efraín Díaz}
\label{transcripcion:ent02}

Modalidad: presencial. Duración aproximada: 20 minutos.\par\smallskip

\noindent\textbf{Entrevistadores (0:} 07): ¿Nos permite grabar la sesión?\par\smallskip

\noindent\textbf{Efraín Díaz (0:} 13): Ok, listo.\par\smallskip

\noindent\textbf{Entrevistadores (0:} 33): Muy buenas noches. Estamos realizando una entrevista para la materia de Interacción Hombre-Máquina en Ingeniería de Requerimientos, enfocada en el diseño de un sistema de coordinación de sesiones académicas. La plataforma busca optimizar la coordinación entre docentes y estudiantes en encuentros tanto personalizados como grupales.\par\smallskip

\noindent\textbf{Entrevistadores (1:} 26): ¿Qué experiencias ha tenido con la coordinación de sesiones académicas con los estudiantes?\par\smallskip

\noindent\textbf{Efraín Díaz (1:} 37): Lo que tiene que ver con refuerzos generalmente se lleva de forma virtual, porque es un poco complejo trabajarlo de forma presencial por el tema de la disponibilidad de los laboratorios.\par\smallskip

\noindent\textbf{Entrevistadores (2:} 07): Si la disponibilidad de los laboratorios se automatizara, ¿estaría abierto a usar la modalidad presencial?\par\smallskip

\noindent\textbf{Efraín Díaz (2:} 20): Sí, por supuesto, que es lo mejor.\par\smallskip

\noindent\textbf{Entrevistadores (2:} 25): ¿Con qué frecuencia realiza sesiones académicas con los estudiantes?\par\smallskip

\noindent\textbf{Efraín Díaz (2:} 36): En todas las materias que doy — programación y administración de bases de datos. La frecuencia depende del tema y del tipo de revisión: reforzar conocimientos o retroalimentar sobre una evaluación.\par\smallskip

\noindent\textbf{Entrevistadores (3:} 50): ¿Considera adecuado el correo institucional como método de acceso, o qué otra forma de autenticación recomendaría?\par\smallskip

\noindent\textbf{Efraín Díaz (4:} 01): Lo ideal sería acceso único con el usuario y clave del SGA, porque valida más aspectos y podría obtener información adicional del curso, paralelo y materia del estudiante.\par\smallskip

\noindent\textbf{Entrevistadores (8:} 41): ¿Considera necesario que el sistema permita aceptar, rechazar o reprogramar solicitudes?\par\smallskip

\noindent\textbf{Efraín Díaz (8:} 56): Sí, en la parte operativa.\par\smallskip

\noindent\textbf{Entrevistadores (9:} 14): ¿El sistema debería permitir a los docentes registrar un horario para que los estudiantes seleccionen?\par\smallskip

\noindent\textbf{Efraín Díaz (9:} 36): Sí, pero es relativo. Los docentes tenemos horas complementarias, pero no siempre están libres. Sería conveniente que cada profesor especifique un grupo de horas disponibles fuera de la jornada de 40 horas.\par\smallskip

\noindent\textbf{Entrevistadores (10:} 46): ¿Cuál canal de comunicación considera más útil para recibir notificaciones?\par\smallskip

\noindent\textbf{Efraín Díaz (11:} 07): El más directo es WhatsApp. Un correo no lo leo de inmediato. Una aplicación móvil con alertas también sería útil.\par\smallskip

\noindent\textbf{Entrevistadores (11:} 45): ¿Qué tan útil sería un recordatorio automático antes de la sesión?\par\smallskip

\noindent\textbf{Efraín Díaz (12:} 04): Debería ser configurable. No solo un día antes, sino que cada usuario pueda definir si quiere 10, 15 minutos antes o más.\par\smallskip

\noindent\textbf{Entrevistadores (12:} 39): ¿Las sesiones deberían tener duración configurable?\par\smallskip

\noindent\textbf{Efraín Díaz (12:} 53): Sí, depende del tema. Puede durar 10 minutos o hasta tres horas.\par\smallskip

\noindent\textbf{Entrevistadores (13:} 19): ¿Las sesiones deberían ser personalizadas o grupales?\par\smallskip

\noindent\textbf{Efraín Díaz (13:} 33): Depende. Lo ideal sería grupal para aprovechar recursos, pero también debe permitirse lo individual.\par\smallskip

\noindent\textbf{Entrevistadores (14:} 04): ¿Qué tan útiles serían reportes sobre las sesiones académicas?\par\smallskip

\noindent\textbf{Efraín Díaz (14:} 13): Serían importantes como evidencia del trabajo extra realizado.\par\smallskip

\noindent\textbf{Entrevistadores (15:} 18): ¿Considera necesario que la plataforma genere reportes automáticos en PDF o Excel?\par\smallskip

\noindent\textbf{Efraín Díaz (15:} 26): Sí, en PDF.\par\smallskip

\noindent\textbf{Entrevistadores (15:} 32): ¿Qué plazo sería razonable para responder a una solicitud?\par\smallskip

\noindent\textbf{Efraín Díaz (15:} 43): Un plazo de 24 horas. No podría ser inmediato si estoy dando clases.\par\smallskip

\noindent\textbf{Entrevistadores (16:} 06): ¿Qué opina sobre limitar el número de solicitudes que un estudiante puede hacer?\par\smallskip

\noindent\textbf{Efraín Díaz (16:} 17): No, debe quedar libre. Eso depende del acuerdo entre docente y estudiante.\par\smallskip

\noindent\textbf{Entrevistadores (16:} 45): ¿Considera importante registrar la asistencia de los estudiantes en las sesiones?\par\smallskip

\noindent\textbf{Efraín Díaz (16:} 54): Sí, o al menos los participantes de cada sesión.\par\smallskip

\noindent\textbf{Entrevistadores (18:} 12): ¿Sería útil que la plataforma permita usar distintas herramientas de videoconferencia como Zoom, Teams o Google Meet?\par\smallskip

\noindent\textbf{Efraín Díaz (18:} 33): Eso es indiferente, depende de lo que tenga cada docente.\par\smallskip

\noindent\textbf{Entrevistadores (18:} 53): ¿Qué nos recomendaría para el desarrollo del proyecto?\par\smallskip

\noindent\textbf{Efraín Díaz (19:} 04): Revisen bien el contexto y el alcance de lo que están haciendo, para que en base a eso definan los elementos, requisitos y avancen con mayor claridad.\par\smallskip

\subsection{ENT-03 — Angelita Bosquez}
\label{transcripcion:ent03}

Modalidad: virtual. Docente de Álgebra y Análisis de Funciones.\par\smallskip

\noindent\textbf{Entrevistadores:} ¿Ha recibido solicitudes de tutoría por parte de los estudiantes?\par\smallskip

\noindent\textbf{Docente:} Sí, he recibido solicitudes. Con ellos tengo un encuentro semanal a través de una clase síncrona. No me envían solicitudes formales, ya que saben que las tutorías están planificadas cada semana por videollamada.\par\smallskip

\noindent\textbf{Entrevistadores:} ¿Cuál fue el método que utilizó para coordinar las tutorías?\par\smallskip

\noindent\textbf{Docente:} Como al inicio no conocía a los estudiantes, cargué a través de la SGEA a los 64 que estaban bajo mi responsabilidad. Luego formé un grupo de WhatsApp. Una vez elegida la presidenta del curso, me retiré del grupo y desde entonces me comunico directamente con ella, salvo situaciones de contacto inmediato, en las que uso el correo institucional.\par\smallskip

\noindent\textbf{Entrevistadores:} ¿Qué dificultades enfrentó al principio para coordinar las tutorías?\par\smallskip

\noindent\textbf{Docente:} Principalmente verificar que los números de teléfono proporcionados fueran correctos, ya que muchos no lo eran. Eso dificultó la comunicación inicial.\par\smallskip

\noindent\textbf{Entrevistadores:} ¿Cómo maneja actualmente la disponibilidad de horarios para las tutorías?\par\smallskip

\noindent\textbf{Docente:} Tenemos los miércoles establecidos en el horario para las tutorías. Desde el inicio se acordó ese día, por lo que no hay inconvenientes. Sin embargo, no siempre se conectan todos los estudiantes, porque muchos trabajan.\par\smallskip

\noindent\textbf{Entrevistadores:} ¿Considera que sería útil automatizar este proceso?\par\smallskip

\noindent\textbf{Docente:} Claro que sí, porque nos ahorraría tiempo y esfuerzo.\par\smallskip

\noindent\textbf{Entrevistadores:} ¿Qué funciones considera más importantes para un sistema de gestión de tutorías?\par\smallskip

\noindent\textbf{Docente:} Que me alerte si un estudiante no está cumpliendo con sus actividades y que también me notifique sobre la asistencia. Aunque sea una carrera en línea, la asistencia sigue siendo importante.\par\smallskip

\noindent\textbf{Entrevistadores:} ¿Cómo le gustaría recibir notificaciones de solicitudes de tutoría?\par\smallskip

\noindent\textbf{Docente:} A través del correo electrónico, ya que es un medio más formal.\par\smallskip

\noindent\textbf{Entrevistadores:} ¿Qué tipo de reportes le gustaría visualizar en el sistema?\par\smallskip

\noindent\textbf{Docente:} Reportes por asignatura, académicos, de notas y de asistencia. Además, que reflejen alertas visuales (amarillo, rojo) dependiendo del porcentaje de cumplimiento.\par\smallskip

\noindent\textbf{Entrevistadores:} ¿Qué reglas o límites considera necesarios en las tutorías?\par\smallskip

\noindent\textbf{Docente:} Que el enlace de acceso no sea compartido con personas externas, de modo que solo participen los estudiantes correspondientes.\par\smallskip

\noindent\textbf{Entrevistadores:} ¿Cuál cree que sería el impacto de un sistema de gestión de tutorías en la labor docente?\par\smallskip

\noindent\textbf{Docente:} Sería de gran impacto, porque permitiría reducir el esfuerzo humano, optimizar el tiempo y dar respuestas más rápidas, facilitando así el proceso de tutoría y el apoyo a los estudiantes.\par\smallskip

\subsection{ENT-05 — Carlos Márquez de la Plata}
\label{transcripcion:ent05}

Impartía cuatro asignaturas: Ventilación de Operaciones A y B, Metodológicos A y B.\par\smallskip

\noindent\textbf{Entrevistadores (0:} 49): ¿Ha recibido solicitudes de tutoría por parte de los estudiantes?\par\smallskip

\noindent\textbf{Docente (0:} 53): Solicitudes como tal, no.\par\smallskip

\noindent\textbf{Entrevistadores (0:} 58): Sabiendo eso, ¿qué método utilizaría para coordinar las tutorías?\par\smallskip

\noindent\textbf{Docente (1:} 06): Si un estudiante tiene alguna dificultad, podemos hacer un grupo y trabajar una tutoría virtual.\par\smallskip

\noindent\textbf{Entrevistadores (1:} 25): ¿Qué dificultades cree que enfrentaría al coordinar las tutorías?\par\smallskip

\noindent\textbf{Docente (1:} 32): Basado en la experiencia del COVID, a veces los estudiantes no prendían las cámaras y no se podía verificar su presencia real.\par\smallskip

\noindent\textbf{Entrevistadores (1:} 54): ¿Cómo manejaría su disponibilidad de horario?\par\smallskip

\noindent\textbf{Docente (1:} 59): Si hay un grupo de estudiantes, se ponen de acuerdo. Con tal que no haya choque de horario y no tenga clases, no importa.\par\smallskip

\noindent\textbf{Entrevistadores (2:} 14): ¿Le gustaría que el proceso de tutoría sea automatizado en vez de coordinar todo manualmente?\par\smallskip

\noindent\textbf{Docente (2:} 24): Sería magnífico, sí.\par\smallskip

\noindent\textbf{Entrevistadores (2:} 28): ¿Cuáles cree que serían las funciones más importantes que debería tener el sistema?\par\smallskip

\noindent\textbf{Docente (2:} 37): El estado de los estudiantes, la asignación de hora previa y el medio de comunicación.\par\smallskip

\noindent\textbf{Entrevistadores (3:} 17): ¿Qué tipo de reportes le gustaría ver como docente?\par\smallskip

\noindent\textbf{Docente (3:} 41): La asistencia a la tutoría, los temas recurrentes que solicitan más los estudiantes, y también un cuadro de diferencia de antes y después de recibir la tutoría, para ver cuánto mejora el estudiante. Un plan de diagnóstico también podría ser.\par\smallskip

\noindent\textbf{Entrevistadores (4:} 15): ¿Qué reglas o límites cree que deberían establecerse en las tutorías?\par\smallskip

\noindent\textbf{Docente (4:} 20): La puntualidad y el comportamiento de la persona.\par\smallskip

\noindent\textbf{Entrevistadores (4:} 29): ¿La tutoría debe ser de forma individual o grupal?\par\smallskip

\noindent\textbf{Docente (4:} 39): Prefiero que sea grupal — puede ser tres o cuatro personas que estén interesadas. Pero también debería permitirse lo individual para el estudiante tímido.\par\smallskip

\noindent\textbf{Entrevistadores (5:} 18): ¿Cuál cree que sería el impacto del sistema propuesto en la labor docente?\par\smallskip

\noindent\textbf{Docente (5:} 24): Sería muy útil. Si se aplica correctamente, sería muy útil.\par\smallskip

\subsection{ENT-04 — Daisy Nata}
\label{transcripcion:ent04}

Asignaturas: Legislación Informática, Estructuras Discretas II, Fundamentos de Física para la Ingeniería, Ciencia Tecnología e Innovación.\par\smallskip

\noindent\textbf{Entrevistadores:} ¿Ha recibido solicitudes de tutoría por parte de los estudiantes?\par\smallskip

\noindent\textbf{Docente:} No he recibido solicitudes formales de tutorías. Lo que hacemos normalmente es que, cuando un tema no queda claro en clase, se lo vuelve a explicar y se refuerza con ejercicios. Solicitudes de tutorías como tal no se me han presentado, y actualmente no tengo asignado ningún grupo como tutor académico.\par\smallskip

\noindent\textbf{Entrevistadores:} Si llegara a recibir solicitudes, ¿qué medios utilizaría para coordinarlas?\par\smallskip

\noindent\textbf{Docente:} Principalmente el correo electrónico, ya que en el SGA actualmente no existe una opción específica para que los estudiantes soliciten tutorías. Por lo general convoco a los estudiantes indicando el tema y la fecha.\par\smallskip

\noindent\textbf{Entrevistadores:} ¿Qué dificultades cree que enfrentaría al coordinar estas tutorías?\par\smallskip

\noindent\textbf{Docente:} La principal dificultad sería el horario. Muchas veces los estudiantes tienen otras clases o actividades, lo que genera choques. Lo ideal sería que el sistema pueda integrarse con el SGA y verificar automáticamente los horarios, evitando conflictos.\par\smallskip

\noindent\textbf{Entrevistadores:} ¿Cómo maneja actualmente su disponibilidad para refuerzos académicos?\par\smallskip

\noindent\textbf{Docente:} Con los estudiantes de la mañana, suelo programar tutorías en la tarde, dentro de mis horas disponibles de preparación de clases. Sin embargo, también debo cumplir con actividades administrativas y de gestión, lo que limita los espacios.\par\smallskip

\noindent\textbf{Entrevistadores:} ¿Le parecería útil automatizar este proceso?\par\smallskip

\noindent\textbf{Docente:} Sí, sería muy útil. Podría contemplarse tanto para tutorías presenciales como virtuales, y permitiría llevar un control del cumplimiento tanto del estudiante como del docente.\par\smallskip

\noindent\textbf{Entrevistadores:} ¿Qué funcionalidades debería tener un sistema de gestión de tutorías?\par\smallskip

\noindent\textbf{Docente:} Lo más importante es que se complemente con el SGA, de manera que pueda gestionar los horarios de los estudiantes y docentes, evitando convocatorias en horarios ya ocupados. También debería generar reportes de cumplimiento de tutorías, asistencia y participación.\par\smallskip

\noindent\textbf{Entrevistadores:} ¿Por qué medio le gustaría recibir notificaciones de solicitudes?\par\smallskip

\noindent\textbf{Docente:} Por correo institucional, pero también a través de notificaciones dentro del SGA, como ocurre actualmente con las tareas.\par\smallskip

\noindent\textbf{Entrevistadores:} ¿Qué tipo de reportes le gustaría visualizar en el sistema?\par\smallskip

\noindent\textbf{Docente:} Reportes que reflejen tanto el trabajo del docente como el cumplimiento de los estudiantes — porcentaje de asistencia y participación en tutorías. Esto permitiría evidenciar si los llamados a tutorías están siendo efectivos o no.\par\smallskip

\noindent\textbf{Entrevistadores:} ¿Qué reglas o límites considera necesarios?\par\smallskip

\noindent\textbf{Docente:} Normas de asistencia mínima. Por ejemplo, si convoco a tres estudiantes y solo uno asiste de manera constante, se reflejaría en el porcentaje de avance. Esto permitiría un seguimiento justo y equilibrado.\par\smallskip

\noindent\textbf{Entrevistadores:} ¿Qué impacto cree que tendría un sistema de este tipo en la labor docente?\par\smallskip

\noindent\textbf{Docente:} Sería de gran impacto porque permitiría llevar un control organizado de las tutorías y reforzaría los aprendizajes por asignatura. Además, sería una innovación impulsada desde la carrera de Software que luego podría replicarse en otras carreras de la universidad.\par\smallskip

\noindent\textbf{Entrevistadores:} Muchas gracias, ingeniera, por su tiempo y colaboración.\par\smallskip

\noindent\textbf{Docente:} Con gusto. Que tengan un buen día.\par\smallskip

\endgroup
